% =====================================================================
% TEMPLATE: TCC / MONOGRAFIA (Unespar — Licenciatura em Matemática)
% Baseado em: projects/[TCC] - Geometria Riemanniana
% Compilar com: latexmk -xelatex -outdir=../build main.tex
% =====================================================================
\documentclass[12pt, oneside]{book}

% =====================================================
% Preâmbulo (pacotes, ambientes de teorema, etc.)
% =====================================================
% !TEX root = ./main.tex
% =====================================================
% PREÂMBULO — TCC
% Os pacotes vêm de comum/ (ver comum/packages.tex).
% Coloque aqui só o que for específico deste trabalho.
% =====================================================

% ---------- Configuração dos módulos comuns ----------
\newcommand{\comumNumeracao}{chapter}        % Teorema 2.1, Definição 2.3, ...
\newcommand{\comumBibliografia}{referencias.bib}
% \newcommand{\comumContadorUnico}{}         % um só contador para tudo
% \newcommand{\comumAsymptote}{}             % liga o pacote asymptote

\input{comum/packages}

% ---------- Específico deste trabalho ----------
% Ex.: \usetikzlibrary{...}, \DeclareMathOperator{\Hom}{Hom}, ...


% =====================================================
% Dados do trabalho (usados na capa e nos metadados do PDF)
% =====================================================
\newcommand{\tituloTCC}{Título do Trabalho de Conclusão de Curso}
\newcommand{\autorUm}{Nome do Autor}
\newcommand{\autorDois}{Nome do Orientador}
\newcommand{\autorTres}{Nome do Coorientador}

\title{\tituloTCC}
\author{\autorUm \and \autorDois \and \autorTres}
\date{\today}

\hypersetup{
  pdftitle={\tituloTCC},
  pdfauthor={\autorUm, \autorDois, \autorTres}
}

\begin{document}

% =====================================================
% Elementos pré-textuais
% =====================================================
% !TEX root = ../main.tex

% =====================================================
% CAPA
% =====================================================
\begin{titlepage}
  \centering
  {\large UNIVERSIDADE ESTADUAL DO PARANÁ\par}
  {\large CURSO DE LICENCIATURA EM MATEMÁTICA\par}
  \vspace{4cm}

  {\Huge \bfseries \tituloTCC \par}
  \vspace{3cm}

  {\Large
    \autorUm \\[0.3cm]
    \autorDois \\[0.3cm]
    \autorTres \par
  }

  \vfill

  {\large
    Trabalho de Conclusão de Curso apresentado como \\
    requisito parcial para obtenção do título de \\
    Licenciado em Matemática.\par
  }

  \vspace{2cm}
  {\large Campo Mourão \\ \the\year \par}
\end{titlepage}


\frontmatter
\tableofcontents
\clearpage

% (opcional) Resumo/Abstract
% \input{capitulos/resumo}

% =====================================================
% Corpo do texto — cada capítulo em um arquivo próprio
% =====================================================
\mainmatter

% !TEX root = ../main.tex

\chapter{Introdução}
\label{cap:introducao}

Texto da introdução. Cite com \verb|\cite|, \verb|\textcite| ou
\verb|\parencite|: segundo \textcite{doCarmo1992}, \ldots{} \parencite{lima1973variedades}.

Referências cruzadas usam \verb|\cref|/\verb|\Cref|: veja a
\cref{def:exemplo} e o \Cref{teo:exemplo}.

\section{Exemplos de ambientes}

\begin{definition}
  \label{def:exemplo}
  Um \emph{objeto} é \ldots
\end{definition}

\begin{theorem}
  \label{teo:exemplo}
  Todo objeto satisfaz \ldots
\end{theorem}

\begin{proof}
  Demonstração.
\end{proof}

\begin{observation}
  Uma observação.
\end{observation}

\begin{citacao}
  Citação longa (mais de três linhas) no padrão ABNT: recuo de 4\,cm,
  fonte menor e espaçamento simples \parencite[p.~1]{doCarmo1992}.
\end{citacao}

% \input{capitulos/cap2_preliminares}
% ...
% !TEX root = ../main.tex

\chapter{Considerações finais}
\label{cap:conclusao}

Texto das considerações finais.


% =====================================================
% Referências bibliográficas
% =====================================================
\printbibliography[
  title={Referências Bibliográficas},
  heading=bibintoc
]
\end{document}
